%generelle Dopplung:
\subsection{Zielsetzung}

Vor diesem Hintergrund verfolgt die vorliegende Arbeit das Ziel, die Implementierung und den Einsatz eines \ac{KI}-Agenten in dem \ac{ERP}-System \ac{D365FnSCM} im Hinblick auf Effizienzpotenziale zu analysieren und zu bewerten, um zu einer Beantwortung der Hauptforschungsfrage zu gelangen.

Die Arbeit verfolgt dabei nicht das Ziel, ein produktiv einsetzbares System zu entwickeln, sondern einen prototypischen Implementierungsansatz zu untersuchen. Die gewonnenen Ergebnisse sollen exemplarische Erkenntnisse zum Einsatz von \ac{KI}-Agenten in \ac{ERP}-Systemen liefern.

Um diese vollumfänglich beantworten zu können, werden zunächst die beiden Unterfragen betrachtet. Hierfür wird zur Identifikation der technischen und organisatorischen Herausforderungen sowie Chancen bei der Implementierung von \ac{KI}-Agenten in \ac{ERP}-Systemen eine Literaturrecherche vorgenommen. Ergänzend dazu werden aus durchgeführten Experteninterviews weitere Erkenntnisse gewonnen, um einen umfassenden Überblick zu erhalten. Darauf aufbauend wird ein \ac{KI}-Agent für ein konkretes Anwendungsszenario in dem \ac{ERP}-System \ac{D365FnSCM} mit Hilfe von Copilot Studio entwickelt und implementiert. 